% Compile: pdflatex monitor-tool.tex && pdflatex monitor-tool.tex

\documentclass[11pt]{article}

\usepackage[margin=1in]{geometry}
\usepackage[utf8]{inputenc}
\usepackage[T1]{fontenc}
\usepackage{hyperref}
\usepackage{url}
\usepackage{booktabs}
\usepackage{amsmath}
\usepackage{amsfonts}
\usepackage{microtype}
\usepackage{enumitem}
\usepackage{graphicx}
\usepackage{xcolor}

\hypersetup{
  colorlinks=true,
  linkcolor=blue!60!black,
  citecolor=blue!60!black,
  urlcolor=blue!60!black,
}

\title{The Monitor Tool: Supervising Long-Running Processes}

\author{
  XSEC Labs \\
  \texttt{github.com/uncesaii/xsec}
}

\date{}

\begin{document}

\maketitle

\begin{abstract}
XSEC's agent can run commands, but it cannot \textit{supervise} a process that
outlives a single tool call. Every execution primitive shipped today is built
around ``run to completion within a timeout, then hand back the whole output,'' a
model that breaks the moment the agent must stand up a dev server and then attack
it, drive a long scanner that emits findings incrementally, watch a fuzzer for
the first crash marker, or kick off a build and continue working. This paper
specifies a dedicated \texttt{monitor} tool that makes long-running supervision a
first-class, structured capability. We audit the current execution surface and
show that a background-run + poll/tail + wait-until primitive is net-new; we
study prior art (Oh My Pi's \texttt{hub} process family, whose structured
ready-gate is the model to follow, and Claude Code's Monitor tool, whose lack of
structured readiness is the instructive weakness); and we propose a single
namespaced \texttt{monitor} tool with an \texttt{op} discriminator and three
deliberate tightenings for a security agent: \texttt{args[]} instead of a shell
string, a structured validated ready-gate, and signal-first cursor-bounded
reporting. Every repository claim is anchored to a
\texttt{packages/<path>:<line>} reference.
\end{abstract}

\section{Problem}

XSEC's agent can run commands, but it cannot \textit{supervise} a process that
outlives a single tool call. Every execution primitive we ship today is built
around ``run to completion within a timeout, then hand back the whole output.''
That model breaks the moment the agent needs to:

\begin{itemize}[leftmargin=*]
  \item Stand up a \textbf{dev server / target under test} and then attack it
  (the server must stay up across many subsequent tool calls).
  \item Drive a \textbf{long scanner} --- \texttt{nuclei}, \texttt{nmap},
  \texttt{ffuf}, \texttt{gobuster}, \texttt{sqlmap} --- that legitimately runs for
  minutes and emits findings incrementally.
  \item Run a \textbf{fuzzer} and watch for the first crash/finding marker, then
  stop it.
  \item Kick off a \textbf{build} or install step and continue working while it
  compiles, coming back only when it finishes (or fails).
\end{itemize}

Right now the agent fakes all of this with \texttt{bash} timeouts and shell
plumbing. This doc specifies a dedicated \texttt{monitor} tool that makes
long-running supervision a first-class, structured capability.

\section{The gap in the current tool surface}

We audited every execution-capable tool in the engine. None of them background a
process and let the agent poll, tail, or wait on it.

\begin{itemize}[leftmargin=*]
  \item \textbf{\texttt{bash}}
  (\texttt{packages/core/src/agent/tools/system.ts:104}) runs a shell command
  with a \textbf{timeout only} --- default 30s, max 120s
  (\texttt{system.ts:110}). It spawns the command in a \textbf{detached process
  group purely so the timeout can kill the whole tree}, not so it can be left
  running. There is no handle, no polling, no tail. When the call returns, the
  process is dead.
  \item \textbf{\texttt{run\_command}} (registered at
  \texttt{packages/core/src/agent/tools/index.ts:79}, defined in
  \texttt{system.ts:78}) is the read-only analysis variant: a restricted command
  surface (``No shell operators like \texttt{;}, \texttt{\&\&}, \texttt{<},
  \texttt{>}, \texttt{\$}'', \texttt{system.ts:80}) with the same
  run-to-completion-within-a-timeout contract. Also not a supervisor.
  \item \textbf{\texttt{pty\_session}} (\texttt{index.ts:96}) is the closest
  existing primitive: a \textit{persistent interactive terminal} with create /
  send / read / close / list operations. But it is \textbf{interactive-shell
  oriented} --- it models a human at a TTY, not a managed daemon with readiness
  and lifecycle --- and it is gated behind the \texttt{ptySession} feature flag,
  \textbf{default OFF}.
  \item \textbf{\texttt{python\_exec}} (\texttt{index.ts:117}) is a persistent
  kernel, but it is \textbf{compute-only} (a Python REPL state), not a process
  supervisor, and is also \textbf{default OFF}.
  \item \textbf{\texttt{spawn\_persistent\_agent}} (\texttt{index.ts:88}) is an
  \textbf{agent} lifecycle --- a long-lived \textit{reasoning} worker --- not an
  OS-process monitor.
\end{itemize}

So a background-run + poll/tail + wait-until primitive is \textbf{net-new}.
Nothing in the registry composes into it without gluing shell hacks together
inside \texttt{bash}.

\section{Prior art}

\subsection{The model to follow --- Oh My Pi's \texttt{hub} process family}

Oh My Pi ships exactly the primitive we want, and its shape is worth copying
closely. See the \texttt{hub} tool docs \citep{omp-hub}. The \texttt{hub} process
family exposes: \texttt{start}, \texttt{ps}, \texttt{logs}, \texttt{stop},
\texttt{restart}, \texttt{describe}, \texttt{send}, \texttt{wait}.

The standout primitive is the \textbf{ready-gate on \texttt{start}}. Instead of
returning as soon as the process is spawned, \texttt{start} accepts a
\texttt{ready} condition and returns \textit{only once that condition is met}:

\begin{itemize}[leftmargin=*]
  \item \texttt{ready.log} --- a \textbf{regex} matched against the process's
  output.
  \item \texttt{ready.port} --- a TCP port (1--65535) that must accept a
  connection.
  \item \texttt{ready.host} --- host for the port probe.
  \item \texttt{ready.timeout} --- default 30s, capped at 1h.
\end{itemize}

If both \texttt{log} and \texttt{port} are supplied, the gate is an \textbf{AND}:
the process is ``ready'' only when the log line has appeared \textit{and} the port
is accepting. This is the difference between ``I launched a server'' and ``the
server is actually serving.''

Its \texttt{logs} operation is equally well-shaped:

\begin{itemize}[leftmargin=*]
  \item \texttt{grep} --- regex filter over the buffered output.
  \item \texttt{lines} --- tail N lines.
  \item \texttt{follow} --- stream until the log advances or the process exits.
  \item \texttt{cursor} --- a \textbf{byte offset to resume from}, so repeated
  polls never re-emit output the agent has already seen. Each call returns a
  \textbf{new cursor}.
  \item \texttt{timeout} --- default 30s, max 3600s.
\end{itemize}

\texttt{send} writes to a running process: \texttt{text} (to stdin), \texttt{keys}
(named keys like \texttt{ENTER} / \texttt{CTRL\_C}), or a \texttt{signal}.

\texttt{wait} \textbf{races} three outcomes --- process exit, an output-pattern
regex match, and a timeout --- and also supports a \textbf{wait-for-all} mode
across several processes. Critically, \texttt{wait} must \textbf{track process
identity across auto-restart}: if a supervised process crashes and is restarted, a
pending \texttt{wait} on the \textit{old} process must not be spuriously satisfied
by the \textit{replacement}. Identity is per incarnation, not per name.

\subsection{The contrast --- Claude Code's Monitor tool}

Claude Code has a Monitor tool \citep{cc-monitor} that runs a shell command in the
background and \textbf{streams stdout as event notifications}, with lines batched
within a $\sim$200ms window. Its parameters are \texttt{command},
\texttt{description}, \texttt{timeout\_ms} (default 300000, cap 3600000), and
\texttt{persistent}.

The instructive weakness: Monitor has \textbf{no structured readiness or wait
condition}. To express ``wait until the server prints \texttt{Listening}'' you
encode a \texttt{while} loop \textit{inside the command string} --- polling
\texttt{curl} or grepping a log in shell. That works, but it pushes the control
logic into a fragile, injection-prone shell one-liner that the agent has to author
correctly every time. Claude Code's background \texttt{Bash}
(\texttt{run\_in\_background}) is even thinner: detach and notify \textbf{once} on
completion.

The lesson from both: streaming and backgrounding are necessary but not
sufficient. The \textbf{structured ready-gate and structured wait} are what turn
``a process is running'' into ``a condition the agent can reason about.''

\section{Proposed spec}

A single namespaced \texttt{monitor} tool with an \texttt{op} discriminator,
mirroring \texttt{hub} but tightened for a security tool. Operations:
\texttt{start}, \texttt{logs}, \texttt{wait}, \texttt{stop}, \texttt{restart},
\texttt{ps}, \texttt{send}.

\subsection{\texttt{start}}

Launch and supervise a process, blocking until ready (or the ready-gate times
out).

\begin{table}[ht]
\centering
\small
\begin{tabular}{lll}
\toprule
field & type & notes \\
\midrule
\texttt{name} & string, 1--48 chars & stable handle used by every other op \\
\texttt{application} & string & the executable --- \textbf{not} a shell string \\
\texttt{args} & string[] & argv vector, passed directly to exec \\
\texttt{env} & object & extra environment variables \\
\texttt{cwd} & string & working directory \\
\texttt{ready.log} & regex & output line that marks readiness \\
\texttt{ready.port} & 1--65535 & TCP port that must accept \\
\texttt{ready.host} & string & host for the port probe \\
\texttt{ready.timeout\_s} & number & default 30, cap 3600 \\
\texttt{restart} & enum & \texttt{no} \textbar\ \texttt{on-failure} \textbar\ \texttt{always} \\
\bottomrule
\end{tabular}
\caption{\texttt{start} operation fields.}
\label{tab:start}
\end{table}

\texttt{start} blocks until the ready condition is satisfied. When both
\texttt{ready.log} and \texttt{ready.port} are given, semantics are \textbf{AND}
--- both must hold. No \texttt{command} shell string is ever accepted (see
rationale below).

\subsection{\texttt{logs}}

Pull filtered, cursor-bounded output from a supervised process.

\begin{table}[ht]
\centering
\small
\begin{tabular}{lll}
\toprule
field & type & notes \\
\midrule
\texttt{name} & string & which process \\
\texttt{grep} & regex & filter output lines \\
\texttt{tail} / \texttt{lines} & number & last N lines \\
\texttt{follow} & bool & stream until advancement or exit \\
\texttt{cursor} & number & byte offset to resume from (no re-emit) \\
\texttt{timeout\_s} & number & default 30, cap 3600 \\
\bottomrule
\end{tabular}
\caption{\texttt{logs} operation fields.}
\label{tab:logs}
\end{table}

Returns: filtered lines \textbf{plus a new cursor}.

\subsection{\texttt{wait}}

Race process outcomes until one resolves.

\begin{table}[ht]
\centering
\small
\begin{tabular}{lll}
\toprule
field & type & notes \\
\midrule
\texttt{name} / \texttt{names} & string / string[] & one or many processes \\
\texttt{for} & enum & \texttt{exit} \textbar\ \texttt{ready} \\
\texttt{pattern} & regex & output pattern to wait for \\
\texttt{timeout\_s} & number & cap 3600 \\
\texttt{mode} & enum & \texttt{any} \textbar\ \texttt{all} \\
\bottomrule
\end{tabular}
\caption{\texttt{wait} operation fields.}
\label{tab:wait}
\end{table}

\texttt{wait} resolves on the first of: matching process \textbf{exit},
\texttt{pattern} match in output, or \texttt{timeout}. \texttt{mode: all} requires
every named process to resolve. Process identity is tracked per incarnation, so an
auto-restart never spuriously satisfies a \texttt{wait} bound to the crashed
process.

\subsection{\texttt{send}}

\begin{table}[ht]
\centering
\small
\begin{tabular}{lll}
\toprule
field & type & notes \\
\midrule
\texttt{name} & string & which process \\
\texttt{text} & string & write to stdin \\
\texttt{keys} & string[] & named keys: \texttt{ENTER}, \texttt{CTRL\_C}, \ldots \\
\texttt{signal} & string & POSIX signal \\
\bottomrule
\end{tabular}
\caption{\texttt{send} operation fields.}
\label{tab:send}
\end{table}

\subsection{\texttt{stop} / \texttt{restart} / \texttt{ps}}

\texttt{stop} terminates (graceful signal, then kill). \texttt{restart} cycles a
process, preserving its name and config. \texttt{ps} lists supervised processes
with state, pid/incarnation, uptime, and last exit code.

\section{How it reports back to the model}

Reporting is where a security-tool monitor should differ from a generic one. The
output has to be \textbf{compact, cursor-bounded, and signal-first}.

\begin{itemize}[leftmargin=*]
  \item \textbf{\texttt{start}} returns a compact \textbf{readiness verdict}:
  \texttt{\{ status: ready | timeout | exited, elapsed, matched \}} where
  \texttt{matched} is the log line and/or the port that satisfied the gate. On
  \texttt{exited} it includes the exit code --- a server that dies during boot is
  a common and important case.
  \item \textbf{\texttt{logs} / \texttt{wait}} return \textbf{only} the filtered,
  cursor-bounded output plus a structured status:
  \texttt{\{ status, matched\_pattern, cursor, exit\_code \}}. Output is truncated
  to a \textbf{token budget ($\sim$25K)} with an explicit
  \texttt{"more at cursor: N"} marker so the agent can page rather than blow its
  context on a chatty scanner.
  \item If streaming (\texttt{follow}), lines are \textbf{event-batched
  ($\sim$200ms)} --- the same batching window Claude Code's Monitor uses --- so
  the agent isn't interrupted per line.
  \item \textbf{Security-specific reporting rules:}
  \begin{itemize}[leftmargin=*]
    \item Surface the \textbf{matched ready-regex line / found marker verbatim}.
    For a scanner, the line that matched is very often \textit{the finding itself}
    (\texttt{[high] SQL injection at \ldots}), so it must not be summarized away.
    \item \textbf{Always return the exit code.} A non-zero exit from
    \texttt{nuclei}, \texttt{nmap}, or a fuzzer is signal, not noise --- it
    distinguishes ``scan completed clean'' from ``scan crashed / target killed the
    connection.''
  \end{itemize}
\end{itemize}

\section{Rationale}

\subsection{Why a structured ready-gate beats a bash \texttt{while} loop}

Encoding readiness as a shell loop
(\texttt{until curl -s localhost:8080; do sleep 1; done}) has four failure modes
that a structured gate eliminates:

\begin{enumerate}[leftmargin=*]
  \item \textbf{Correctness is the model's problem, every time.} The agent has to
  re-author the polling logic --- the right endpoint, the right sleep, the right
  exit condition --- on each launch. A structured \texttt{ready} object is
  declarative and validated once.
  \item \textbf{No composite conditions.} ``Log says ready AND port is open'' is
  awkward in shell but a one-line AND in the gate. Servers routinely print
  ``listening'' \textit{before} the socket is actually accepting; the AND is what
  closes that race.
  \item \textbf{No clean verdict.} A \texttt{while} loop that times out just\ldots
  stops. The agent gets ambiguous output and has to infer what happened. The gate
  returns a discrete \texttt{ready | timeout | exited} with the matched evidence.
  \item \textbf{The ``server died during boot'' case is invisible} to a naive
  poll loop --- it keeps polling a dead pid until timeout. The gate watches
  process exit \textit{and} the readiness condition together and reports
  \texttt{exited} with the code immediately.
\end{enumerate}

This is the same design instinct behind the engine's kernel-verify loop: prefer a
\textbf{validated, structured contract} over trusting free-form model output to
get the plumbing right (see \texttt{packages/core/src/agent/CLAUDE.md},
``Structured output --- validate every submission'').

\subsection{Why \texttt{args[]} (no shell) matters for a security tool}

\texttt{start} deliberately takes \texttt{application} + \texttt{args[]} and
\textbf{refuses a shell command string}. This is not ergonomics --- it's a
security boundary:

\begin{itemize}[leftmargin=*]
  \item \textbf{No shell injection surface.} The agent is often building command
  lines from \textit{attacker-influenced data} it just scraped off a target --- a
  discovered hostname, a reflected parameter, a filename from a directory listing.
  If any of that flows into a shell string, a \texttt{; rm -rf} or \texttt{\$(\ldots)}
  in the target's response becomes command execution on \textit{our} host. An
  argv vector is passed straight to \texttt{exec} with no shell interpretation, so
  a malicious value is just an inert argument.
  \item \textbf{Auditability.} \texttt{\{ application: "nuclei", args: ["-u", url,
  "-t", tpl] \}} is trivially inspectable, allowlistable, and loggable per token.
  A shell string is opaque until parsed.
  \item \textbf{Consistency with the existing posture.} \texttt{run\_command}
  already bans shell operators for exactly this reason (\texttt{system.ts:80}).
  The monitor should be \textit{stricter} than \texttt{bash}, not looser, because
  it runs things that live longer and touch the network.
\end{itemize}

\subsection{How it composes with the rest of the engine}

\begin{itemize}[leftmargin=*]
  \item \textbf{\texttt{spawn\_persistent\_agent}} is the reasoning half;
  \texttt{monitor} is the process half. The intended pattern: a persistent agent
  owns a workflow (``attack this app''), uses \texttt{monitor.start} to bring the
  target up with a ready-gate, then uses \texttt{monitor.logs}/\texttt{wait} to
  observe it while it attacks --- the process outlives any single agent turn,
  which is precisely what \texttt{bash} cannot do.
  \item \textbf{The existing scanner wrappers} (\texttt{tools/scanner.ts},
  \texttt{recon.ts}, \texttt{detections.ts}) are run-to-completion abstractions.
  \texttt{monitor} is the substrate for the \textit{long-tail} cases they can't
  cover: a \texttt{ffuf}/\texttt{nuclei} run the agent wants to watch incrementally
  and cut short on first finding, or a fuzzer left running while other work
  proceeds. Over time the noisier wrappers could be re-expressed on top of
  \texttt{monitor} rather than each re-implementing timeout/tail logic.
  \item \textbf{\texttt{pty\_session}} stays the tool for genuinely
  \textit{interactive} TTY work (a \texttt{msfconsole}, an SSH session, anything
  that needs a real terminal). \texttt{monitor} is for \textbf{managed daemons and
  batch jobs} where readiness and exit codes matter more than terminal semantics.
  They are complementary, not overlapping.
\end{itemize}

\section{Integration notes}

\begin{itemize}[leftmargin=*]
  \item \textbf{Zod-validate then reject.} XSEC tools validate their payload
  against an explicit Zod schema and \textbf{reject on mismatch before any side
  effect} --- \texttt{kernel\_run} is the reference implementation
  (\texttt{packages/core/src/agent/tools/kernel-run.ts}; see
  \texttt{kernelRunArgsSchema} / \texttt{validateKernelRunArgs}). The
  \texttt{monitor} schema must be the single source of truth for its payload:
  bound \texttt{name} (1--48), enum \texttt{op}/\texttt{restart}/\texttt{for}/\texttt{mode},
  port range 1--65535, and the \texttt{timeout\_s} caps (30 default / 3600 max).
  \texttt{.strip()} unknown top-level keys --- in particular, reject any
  \texttt{command}/\texttt{shell} key so a model can't smuggle a shell string past
  the \texttt{args[]} contract. A rejection should be fed back as an
  \texttt{is\_error} tool\_result so the model self-corrects, matching the
  kernel-run pattern.
  \item \textbf{Registration.} Add \texttt{"monitor"} to
  \texttt{TOOL\_REGISTRY\_ORDER} (\texttt{packages/core/src/agent/tools/index.ts:66})
  and a corresponding \texttt{DOMAIN\_DEFINITIONS} entry, then a dispatch route in
  \texttt{packages/core/src/agent/tools/dispatch.ts} (per-tool handler, not a
  shared chokepoint --- see the dispatch-table note at \texttt{dispatch.ts:2}).
  Implementation belongs in a new \texttt{tools/monitor.ts} alongside the other
  execution tools.
  \item \textbf{Gate it like \texttt{bash}.} \texttt{monitor} is network- and
  exec-capable --- it spawns arbitrary processes that bind ports and reach the
  network --- so it must sit behind the \textbf{same role/permission gating and
  feature flagging as \texttt{bash}/\texttt{run\_command}}, not be enabled for
  read-only audit roles. Given the persistence and network reach, defaulting the
  flag \textbf{OFF} (as \texttt{pty\_session} and \texttt{python\_exec} are) is the
  conservative starting point.
  \item \textbf{Lifecycle hygiene.} Supervised processes must be tracked in a
  registry that is torn down on scan completion / agent shutdown, so a crashed or
  forgotten dev server never leaks past the run. Enforce a cap on concurrent
  supervised processes.
\end{itemize}

\section{Decision}

Build a single namespaced \texttt{monitor} tool with ops
\texttt{start / logs / wait / stop / restart / ps / send}, modeled on Oh My Pi's
\texttt{hub} family, with three deliberate tightenings for a security agent:

\begin{enumerate}[leftmargin=*]
  \item \textbf{\texttt{args[]}, never a shell string}, to remove the injection
  surface when launching processes from attacker-influenced data.
  \item \textbf{A structured, validated ready-gate} (log-regex AND TCP port,
  capped timeout) instead of agent-authored shell poll loops.
  \item \textbf{Signal-first, cursor-bounded reporting} --- verbatim matched
  lines, always-present exit codes, $\sim$25K-token truncation with ``more at
  cursor.''
\end{enumerate}

Register it in \texttt{TOOL\_REGISTRY\_ORDER} + \texttt{dispatch.ts},
Zod-validate the payload with reject-on-mismatch (kernel-run pattern), and gate it
like \texttt{bash} with the feature flag defaulting OFF.

\section{Open questions}

\begin{itemize}[leftmargin=*]
  \item \textbf{Log retention \& spill-to-disk.} How much per-process output do we
  buffer in memory before spilling to a file, and where does that file live
  relative to the scan workspace? Chatty fuzzers can produce GBs.
  \item \textbf{Concurrent-process cap.} What is the right ceiling, and is it
  per-scan or global? A fan-out of \texttt{spawn\_agents} could each want their own
  supervised target.
  \item \textbf{Restart policy vs. \texttt{wait} identity semantics.} With
  \texttt{restart: always}, a \texttt{wait for: exit} may never resolve. Do we
  surface incarnation count, and does \texttt{wait} observe restarts as events?
  \item \textbf{Port-probe scope.} Do we allow \texttt{ready.host} to be
  non-loopback (probe a remote target's port), or restrict readiness probing to
  localhost-bound processes we launched?
  \item \textbf{Interaction with \texttt{bash} timeouts.} Should \texttt{bash}
  gain a ``hand this off to monitor'' escape hatch when a command exceeds its 120s
  cap (\texttt{system.ts:110}), or do we keep the two surfaces strictly separate?
  \item \textbf{Output redaction.} Since matched lines are surfaced verbatim and
  may contain secrets the scanner extracted, do we apply the engine's existing
  secret-scrubbing before returning \texttt{logs}/\texttt{wait} output?
\end{itemize}

\begin{thebibliography}{99}

\bibitem{omp-hub}
Oh My Pi. \textit{The \texttt{hub} tool}. Project documentation.
\url{https://github.com/can1357/oh-my-pi/blob/main/docs/tools/hub.md}

\bibitem{cc-monitor}
tkou15. \textit{Claude Code Monitor tool} (write-up). Zenn.
\url{https://zenn.dev/tkou15/articles/claude-code-monitor}

\end{thebibliography}

\end{document}
